\documentclass[12pt,a4paper]{article}

\usepackage[utf8]{inputenc}
\usepackage[T1]{fontenc}
\usepackage[indonesian]{babel}
\usepackage{amsmath,amssymb,amsthm,amsfonts}
\usepackage{bm}
\usepackage{booktabs}
\usepackage{array}
\usepackage{longtable}
\usepackage{geometry}
\usepackage{graphicx}
\usepackage{hyperref}
\usepackage{enumitem}
\usepackage{mathtools}

\geometry{margin=2.5cm}

\title{Derivasi PDE Mundur untuk Harga Klaim Derivatif melalui Formula Feynman--Kac}
\author{}
\date{}

\begin{document}

\maketitle

\section*{Ringkasan Eksekutif}

Laporan ini menyajikan derivasi Persamaan Diferensial Parsial (PDE) mundur untuk harga klaim derivatif dengan menggunakan formula Feynman--Kac, sehingga persamaan stokastik dalam file pengguna (Pers.~(2) dan (5)) ditransformasikan menjadi PDE target (secara analogi dengan Pers.~(8) pada referensi lain).

Kita mulai dengan mengekstrak notasi dan persamaan kunci dari file pengguna (termasuk Pers.~(4) yang memuat $f(t,Y,Z)$). Selanjutnya dibangun dinamika SDE dua variabel $(X_t,Y_t)$ sesuai notasi paper, dihitung generator infinitesimal menggunakan Teorema It\^o, dan dituliskan PDE mundur:

\[
\partial_t V+\mathcal L V-rV=0.
\]

Laporan ini memetakan setiap suku dari SDE (2) dan (5) ke dalam suku-suku PDE akhir serta mengidentifikasi fungsi driver:

\[
f(t,Y,Z)=-rY-\theta Z.
\]

Selain itu dibahas beberapa kasus khusus (misalnya $Y$ deterministik versus stokastik, korelasi antar Brownian motion, dan perubahan ukuran risiko), berikut asumsi regularitas yang diperlukan seperti kelas $C^{2,1}$ dan sifat Markov.

Sebagai pelengkap disertakan metode numerik (finite difference, ADI, method of lines, dan Monte Carlo berbasis Feynman--Kac), contoh kerja dengan proses Ornstein--Uhlenbeck, tabel perbandingan asumsi terhadap bentuk PDE, serta diagram alir derivasi.

\section{Persamaan Kunci dalam File Pengguna}

Dari file pengguna (BSDE untuk opsi), diperoleh persamaan berikut.

\subsection*{Persamaan (4)}

Persamaan stokastik diberikan oleh

\begin{equation}
dY(t)=\bigl(rY(t)+\theta Z(t)\bigr)\,dt+Z(t)\,dW(t),
\qquad
Y(0)=y_0.
\tag{4}
\end{equation}

\subsection*{Persamaan (5)--(6)}

Fungsi driver BSDE diberikan oleh

\begin{equation}
f(t,Y(t),Z(t))
=
-rY(t)-\theta Z(t).
\end{equation}

Diasumsikan kondisi terminal

\[
Y(T)=\xi(X(T)).
\]

Persamaan (2) diasumsikan merepresentasikan dinamika aset $S(t)$ dan $B(t)$ yang menghasilkan dinamika portofolio $Y$.

Fokus utama derivasi PDE adalah Pers.~(4) dan identifikasi fungsi driver.

\section{Model SDE dan Generator Infinitesimal}

Definisikan proses stokastik dua dimensi:

\[
(X_t,Y_t)
\]

di bawah ukuran risk-neutral $\mathbb Q$.

\subsection*{Dinamika Variabel $X_t$}

\begin{equation}
dX_t
=
a(X_t,t)\,dt
+
b(X_t,t)\,dW_t^{\mathbb Q},
\end{equation}

dengan:

\begin{itemize}
\item $a(x,t)$ : drift,
\item $b(x,t)$ : volatilitas.
\end{itemize}

Sebagai contoh pada model Ornstein--Uhlenbeck:

\[
a(x,t)=k(\theta-x),
\qquad
b(x,t)=\sigma.
\]

\subsection*{Dinamika Variabel $Y_t$}

\begin{equation}
dY_t
=
f(t,Y_t,Z_t)\,dt
+
Z_t\,dW_t^{\mathbb Q}.
\end{equation}

Dari Pers.~(4):

\[
dY=(rY+\theta Z)\,dt+Z\,dW.
\]

Jika ditulis dalam bentuk BSDE standar:

\[
dY=-f\,dt+Z\,dW,
\]

maka:

\[
f(t,Y,Z)
=
-rY-\theta Z.
\]

\subsection*{Generator Infinitesimal}

Untuk fungsi nilai

\[
V=V(t,x,y),
\]

generator infinitesimalnya adalah

\begin{align}
\mathcal LV
&=
a(x,t)V_x
+
f(t,y,z)V_y
+
\frac12 b(x,t)^2V_{xx}
\nonumber\\
&\quad+
\frac12 c(y,t)^2V_{yy}
+
\rho\,b(x,t)c(y,t)V_{xy}.
\end{align}

Jika $Y$ tidak berdifusi maka:

\[
c(y,t)=0,
\]

sehingga suku:

\[
V_{yy},
\qquad
V_{xy}
\]

menghilang.

\section{Turunan PDE melalui Formula Feynman--Kac}

Menurut teorema Feynman--Kac, PDE mundur diberikan oleh:

\begin{equation}
\frac{\partial V}{\partial t}
+
\mathcal LV
-rV
=
0,
\end{equation}

dengan kondisi terminal

\[
V(T,x,y)=g(x,y).
\]

Substitusi generator menghasilkan:

\begin{align}
\partial_tV
&+
a(x,t)V_x
+
f(t,y,z)V_y
+
\frac12 b(x,t)^2V_{xx}
\nonumber\\
&+
\frac12 c(y,t)^2V_{yy}
+
\rho b(x,t)c(y,t)V_{xy}
-rV
=
0.
\end{align}

Untuk kasus file pengguna diperoleh:

\begin{equation}
\partial_tV
+
a(x,t)V_x
+
(ry+\theta z)V_y
+
\frac12 b(x,t)^2V_{xx}
-rV
=
0.
\end{equation}

dengan kondisi terminal

\[
V(T,x,y)=g(x,y).
\]
\section{Pemetaan Persamaan (2) dan (5) ke PDE Mundur}

Sekarang dicocokkan setiap suku dari Pers.~(2) dan Pers.~(5) pada file pengguna ke bentuk PDE mundur.

\begin{itemize}
\item \textbf{Persamaan (2)} merepresentasikan dinamika aset $X$.

Drift dan volatilitas masuk ke dalam suku:

\[
a(x,t)V_x,
\qquad
\frac12 b(x,t)^2V_{xx}.
\]

Sebagai contoh, apabila model mengikuti proses Ornstein--Uhlenbeck, maka dapat diperoleh:

\[
a(x,t)=k(\theta-x)-\lambda\sigma,
\]

setelah transformasi ke ukuran risk-neutral.

\item \textbf{Persamaan (5)} memberikan driver:

\[
f(t,Y,Z)=-rY-\theta Z.
\]

Karena bentuk BSDE ditulis sebagai

\[
dY=-f\,dt+Z\,dW,
\]

maka drift aktual proses $Y$ menjadi:

\[
\mu_Y=rY+\theta Z.
\]

Akibatnya pada PDE muncul suku:

\[
(ry+\theta z)V_y.
\]

\item \textbf{Diskonto} ditambahkan melalui suku:

\[
-rV.
\]

Suku ini berasal dari prinsip arbitrage-free pricing dalam formula Feynman--Kac.
\end{itemize}

Dengan demikian hubungan antara BSDE dan PDE dapat diringkas sebagai:

\[
f(t,Y,Z)
\longrightarrow
V_y
\]

dan

\[
(rY+\theta Z)
\longrightarrow
(rY+\theta Z)V_y.
\]

\section{Kasus Khusus dan Asumsi}

\subsection*{1. Kasus $Y$ Deterministik}

Jika

\[
dY=f(X,t)\,dt,
\]

maka tidak terdapat komponen stokastik pada $Y$, sehingga:

\[
V_{yy}=0,
\qquad
V_{xy}=0.
\]

Kasus ini identik dengan:

\[
Z=0.
\]

\subsection*{2. Kasus $Y$ Berdifusi}

Jika

\[
dY=f\,dt+c(Y,t)\,dZ,
\]

maka PDE memperoleh tambahan:

\[
\frac12 c(y,t)^2V_{yy}.
\]

\subsection*{3. Korelasi Antar Brownian Motion}

Jika

\[
dW^X\,dW^Y=\rho\,dt,
\]

maka muncul suku campuran:

\[
\rho bcV_{xy}.
\]

\subsection*{4. Perubahan Ukuran Risiko}

Jika model awal dibangun pada ukuran fisik $\mathbb P$, maka drift harus ditransformasikan menjadi:

\[
\tilde a(x)
=
a(x)-\lambda b(x),
\]

sesuai transformasi Girsanov.

\subsection*{5. Asumsi Regularitas}

Diasumsikan:

\begin{enumerate}
\item
\[
V\in C^{2,1},
\]

\item proses bersifat Markov,

\item koefisien kontinu,

\item ekspektasi tetap terbatas.
\end{enumerate}

\section{Metode Numerik dan Verifikasi}

\subsection*{Daftar Periksa Verifikasi}

Pastikan model memenuhi:

\begin{enumerate}
\item Regularitas koefisien:
\[
a,b,f
\]
cukup mulus.

\item Pertumbuhan solusi:
\[
|V|
\]
terkontrol.

\item Kondisi terminal:
\[
V(T)=g(x,y).
\]

\item Boundary condition ditentukan dengan benar.
\end{enumerate}

\subsection*{Metode Penyelesaian}

\begin{enumerate}

\item \textbf{Finite Difference}

Menggunakan skema eksplisit, implisit, atau ADI.

\item \textbf{Method of Lines}

Diskritisasi ruang sehingga PDE menjadi sistem ODE.

\item \textbf{Monte Carlo (Feynman--Kac)}

Simulasi:

\[
(X_t,Y_t)
\]

di bawah ukuran $\mathbb Q$ lalu hitung:

\[
e^{-rT}g(X_T,Y_T).
\]

\item \textbf{Metode Spektral / FEM}

Alternatif untuk dimensi tinggi.

\end{enumerate}

\section{Contoh Kerja: Ornstein--Uhlenbeck dan $Y_t=f(X_t)$}

Misalkan:

\begin{equation}
dX_t
=
k(\theta-X_t)\,dt
+
\sigma\,dW_t.
\end{equation}

dan

\begin{equation}
dY_t=f(X_t)\,dt.
\end{equation}

Generator infinitesimal:

\begin{equation}
\mathcal LV
=
k(\theta-x)V_x
+
f(x)V_y
+
\frac12\sigma^2V_{xx}.
\end{equation}

PDE mundur menjadi:

\begin{equation}
\partial_tV
+
k(\theta-x)V_x
+
f(x)V_y
+
\frac12\sigma^2V_{xx}
-rV
=
0,
\end{equation}

dengan kondisi terminal:

\[
V(T,x,y)=g(x,y).
\]

Langkah derivasi:

\begin{enumerate}

\item Tentukan drift dan difusi:
\[
a(x)=k(\theta-x),
\qquad
b(x)=\sigma.
\]

\item Hitung generator:
\[
\mathcal L.
\]

\item Terapkan Feynman--Kac.

\item Tambahkan kondisi terminal.

\end{enumerate}

\section{Tabel Kasus Asumsi dan Bentuk PDE}

\begin{table}[h]
\centering
\renewcommand{\arraystretch}{1.4}

\begin{tabular}{p{5cm}p{8cm}}

\toprule

\textbf{Asumsi Model}
&
\textbf{Bentuk PDE}

\\
\midrule

$
dY=f(x)\,dt
$

&

$
\partial_tV+aV_x+fV_y+\frac12b^2V_{xx}-rV=0
$

\\

$
dY=fdt+c(y)dZ,\;
\rho=0
$

&

Tambah

$
\frac12c^2V_{yy}
$

\\

$
dY=fdt+c(y)dZ,\;
\rho\neq0
$

&

Tambah

$
\frac12c^2V_{yy}
+
\rho bcV_{xy}
$

\\

Risk-neutral measure

&

$
a\to a-\lambda b
$

\\

Terminal

$
V(T)=g
$

&

Kondisi batas akhir

\\

\bottomrule

\end{tabular}

\end{table}

\section{Diagram Alir Derivasi}

\begin{verbatim}
flowchart TD

A[Mulai:
Tulis SDE]

-->

B[Hitung Generator
dengan Ito]

-->

C[Terapkan
Feynman--Kac]

-->

D[Substitusi
Pers.(2),(5)]

-->

E[Dapatkan PDE]

-->

F[Selesaikan
PDE]
\end{verbatim}

Diagram di atas menggambarkan urutan:

\[
\text{SDE}
\rightarrow
\text{Generator}
\rightarrow
\text{Feynman--Kac}
\rightarrow
\text{PDE}
\rightarrow
\text{Solusi}.
\]

\section*{Referensi}

Teorema Feynman--Kac dan kalkulus stokastik menjadi dasar teoritis derivasi ini.

Referensi yang umum digunakan antara lain:

\begin{enumerate}

\item
B.~{\O}ksendal,
\textit{Stochastic Differential Equations}.

\item
Karatzas dan Shreve,
\textit{Brownian Motion and Stochastic Calculus}.

\item
Shreve,
\textit{Stochastic Calculus for Finance II}.

\end{enumerate}

\end{document}